\begin{longtable}{|c|X|}
    \caption{Special Characters in Documentation Comments} \\
    \hline
    \textbf{Character} & \textbf{Description} \\
    \hline
    \endfirsthead
    \multicolumn{2}{c}%
    {\tablename\ \thetable\ -- \textit{Continued from previous page}} \\
    \hline
    \textbf{Character} & \textbf{Description} \\
    \hline
    \endhead
    \multicolumn{2}{r}{\textit{Continued on next page}} \\
    \endfoot
    \endlastfoot
%------------------------------------------------------------------------------

\makecell{\textbf{\texttt{\textless}} and \textbf{\texttt{\textgreater}}} & 
Used for the HTML tags \\
\hline

%------------------------------------------------------------------------------

\makecell{\textbf{\texttt{\&}}} & 
Introduces an HTML\index{HTML} entity \\
\hline

%------------------------------------------------------------------------------

\makecell{\textbf{\texttt{@}}} & 
The lead-in for a Javadoc\index{javadoc} tag \\
\hline

%------------------------------------------------------------------------------

\makecell{\textbf{\texttt{\}}}} & 
Terminates an inline Javadoc\index{javadoc} tag \\
\hline

%------------------------------------------------------------------------------

\makecell{\textbf{\texttt{"}}} & 
Double-quotes are sometimes problematic for HTML\index{HTML} code \\
\hline

%------------------------------------------------------------------------------

\makecell{\textbf{\texttt{*/}}} & 
This combination terminates the documentation comment \\
\hline

%------------------------------------------------------------------------------

\end{longtable}

%==============================================================================
% End of file!!
%==============================================================================
